\documentclass[UTF8]{ctexart}
\usepackage{amsmath,amssymb}
\usepackage{geometry}
\usepackage{booktabs}
\geometry{a4paper,margin=2.4cm}

\title{序列：期望递推推导}
\author{}
\date{}

\begin{document}

\maketitle

设第 $i$ 次操作成功的概率为
\[
q_i=\frac{p_i}{100}.
\]
再令 $\xi_i$ 为第 $i$ 次操作是否成功的 $0/1$ 随机变量，则
\[
\Pr(\xi_i=1)=q_i,
\qquad
\Pr(\xi_i=0)=1-q_i.
\]
题目中的随机过程可以写成
\[
a_i=\xi_i(a_{i-1}+1).
\]

定义前缀和随机变量
\[
S_i=\sum_{j=1}^{i}a_j.
\]
要求的答案就是
\[
\mathbb E(S_n^2).
\]

\section*{状态定义}

维护以下五个量：
\[
\begin{aligned}
A_i&=\mathbb E(a_i),\\
B_i&=\mathbb E(a_i^2),\\
C_i&=\mathbb E(S_i),\\
D_i&=\mathbb E(S_i a_i),\\
F_i&=\mathbb E(S_i^2).
\end{aligned}
\]

初始时 $a_0=S_0=0$，因此
\[
A_0=B_0=C_0=D_0=F_0=0.
\]

\section*{递推}

由于 $\xi_i$ 与前 $i-1$ 步产生的所有随机变量独立，并且
$\xi_i^2=\xi_i$，所以
\[
\begin{aligned}
A_i
&=\mathbb E\bigl(\xi_i(a_{i-1}+1)\bigr)\\
&=q_i\bigl(\mathbb E(a_{i-1})+1\bigr)\\
&=q_i(A_{i-1}+1).
\end{aligned}
\]

同理，
\[
\begin{aligned}
B_i
&=\mathbb E\bigl(\xi_i^2(a_{i-1}+1)^2\bigr)\\
&=q_i\mathbb E\bigl(a_{i-1}^2+2a_{i-1}+1\bigr)\\
&=q_i(B_{i-1}+2A_{i-1}+1).
\end{aligned}
\]

为了计算 $S_i^2$ 中的交叉项，先计算
\[
H_i=\mathbb E(S_{i-1}a_i).
\]
有
\[
\begin{aligned}
H_i
&=\mathbb E\bigl(S_{i-1}\xi_i(a_{i-1}+1)\bigr)\\
&=q_i\mathbb E\bigl(S_{i-1}a_{i-1}+S_{i-1}\bigr)\\
&=q_i(D_{i-1}+C_{i-1}).
\end{aligned}
\]

因为
\[
S_i=S_{i-1}+a_i,
\]
所以
\[
C_i=C_{i-1}+A_i.
\]

再计算联合期望：
\[
\begin{aligned}
D_i
&=\mathbb E(S_i a_i)\\
&=\mathbb E\bigl((S_{i-1}+a_i)a_i\bigr)\\
&=\mathbb E(S_{i-1}a_i)+\mathbb E(a_i^2)\\
&=H_i+B_i.
\end{aligned}
\]

最后展开平方：
\[
\begin{aligned}
F_i
&=\mathbb E(S_i^2)\\
&=\mathbb E\bigl((S_{i-1}+a_i)^2\bigr)\\
&=\mathbb E(S_{i-1}^2)
  +2\mathbb E(S_{i-1}a_i)
  +\mathbb E(a_i^2)\\
&=F_{i-1}+2H_i+B_i.
\end{aligned}
\]

因此完整递推为
\[
\boxed{
\begin{aligned}
H_i&=q_i(D_{i-1}+C_{i-1}),\\
A_i&=q_i(A_{i-1}+1),\\
B_i&=q_i(B_{i-1}+2A_{i-1}+1),\\
F_i&=F_{i-1}+2H_i+B_i,\\
C_i&=C_{i-1}+A_i,\\
D_i&=H_i+B_i.
\end{aligned}}
\]

\section*{与代码变量的对应}

\begin{center}
\begin{tabular}{cc}
\toprule
代码变量 & 数学含义\\
\midrule
\texttt{p[i]} & $q_i$\\
\texttt{E[i]} & $A_i=\mathbb E(a_i)$\\
\texttt{E2[i]} & $B_i=\mathbb E(a_i^2)$\\
\texttt{sum} & $C_i=\mathbb E(S_i)$\\
\texttt{g} & $D_i=\mathbb E(S_i a_i)$\\
\texttt{x} & $H_i=\mathbb E(S_{i-1}a_i)$\\
\texttt{ans} & $F_i=\mathbb E(S_i^2)$\\
\bottomrule
\end{tabular}
\end{center}

所有递推均在模 $998244353$ 意义下进行。代码用
\[
100^{-1}\pmod {998244353}
\]
把输入的百分数转换为 $q_i$。

最终输出
\[
F_n=\mathbb E\left(\left(\sum_{i=1}^{n}a_i\right)^2\right).
\]

时间复杂度为 $O(n)$。

\end{document}
